% A.116: dyadic proof audited by root and Legendre; odd proof audited by Legendre.
\subsection{An order-four norm calculation on the harmonic eighth-power cover}

\begin{statement}{A.116}{Proved}
{The negative-fourth-power triple character has zero evaluation at every place}
Let $\chi_{16}:G_{\Q}\to\tfrac14\Z/\Z$ be the even
cyclotomic character with $\chi_{16}(5)=1/4$ and
$\chi_{16}(-1)=0$ on $(\Z/16)^\times$.
Let $k=\Q_v$ be any completion of $\Q$.
For $b\in k^\times$ let $\chi_b$ be its quadratic
character, with values $0,1/2$, and put
$\Psi_b=\chi_{16}|_{G_k}+\chi_b$.
Its order divides four; at $v=2$ its order is exactly
four because $2\Psi_b$ is the nonzero quadratic
character of $\Q_2(\sqrt2)$.

Put $F(z)=z(z-1)(z-\tfrac12)$.
Suppose $R,T,c\in k\setminus\{0,1,\tfrac12\}$,
with $c\ne R,T$, and
\[
 \frac TR,\quad\frac{T-1}{R-1},\quad
 \frac{T-\tfrac12}{R-\tfrac12}\in k^{\times8},
 \qquad
                     r^4=-\frac{F(c)}{F(R)}
                       \quad(r\in k^\times).
\tag{A.116.1}
\]
Then the actual cyclic algebra evaluation is zero:
\[
                   \left(\frac{T-c}{R-c},\Psi_r\right)=0
                                    \quad\text{in }\operatorname{Br}(k).
\tag{A.116.2}
\]
For the actual three-root fibre algebra, put
$D_i=(c-c_i)/(R-c_i)$, $c_i=0,1,\tfrac12$.
The product relation is $\prod_iD_i=-r^4$.
The character $\Psi_r$ belongs to the common
abelian field of all factors of
$k[Z_i^{\pm1}]/(Z_i^8-D_i)_i$, by the explicit
product-root map proved below.
This supplies one full order-four character family.
It does not assert that this family generates every
common abelian character of an arbitrary eighth-root fibre.
\end{statement}

\begin{proof}
\emph{The character is in every required root field.}
Every root of $Z^8+b^4$ has the form
$\theta=\zeta_{16}^j\sqrt b$ with $j$ odd.
Write
$g(\zeta_{16})=\zeta_{16}^u$ and
$g(\sqrt b)=(-1)^\epsilon\sqrt b$.
If $g$ fixes this particular root, then
\[
 \begin{gathered}
 j(u-1)+8\epsilon=0\pmod{16},
 \qquad u=1+8\epsilon\pmod{16},\\
            \chi_{16}(g)+\chi_b(g)=
                 \epsilon/2+\epsilon/2=0\pmod{\Z}.
 \end{gathered}
\tag{A.116.3}
\]
Consequently the cyclic field $K_b$ cut out
by $\Psi_b$ is contained in every factor field of
the finite etale algebra $k[Z]/(Z^8+b^4)$.
Its degree is the order of $\Psi_b$, possibly one.

For the three-root fibre in the statement the actual
algebra map is
\[
 k[Z]/(Z^8+r^4)\longrightarrow
       k[Z_i^{\pm1}]/(Z_i^8-D_i)_i,
                     \qquad Z\longmapsto\prod_iZ_i.
\tag{A.116.3a}
\]
Over an algebraic closure its map of finite sets is
the product map from $8^3$ root tuples to the eight
product roots; every fibre has $8^2$ elements.
Thus it is an injective finite etale algebra map of
rank $64$.
Every factor of the target contains the factor of
the source selected by its actual product root.
Taking $b=r$ in (A.116.3) proves that every such target
factor contains $K_r$, which is the asserted
common-character comparison.

It follows by transitivity of the actual field norms that
\[
                       X^8+b^4,\quad 1+b^4
               \quad\hbox{are norms from }K_b/k
\tag{A.116.4}
\]
whenever they are nonzero.
Indeed their values are the products of
$N_{k(\theta)/k}(X-\theta)$ and
$N_{k(\theta)/k}(1-\theta)$ over the factors,
and every one of these norms factors through $K_b$.
For a cyclic extension, a norm makes its cyclic
algebra split: if $a=N(u)$, represent its cyclic
generator by the operator $u\sigma$ on the extension
field, whose $[K_b:k]$th power is multiplication by $a$.
Thus (A.116.4) gives zero for the corresponding symbols.

\emph{The coordinate changes preserve the displayed cyclic class.}
As in A.111 use
$g_1(z)=1-z$ and $g_2(z)=1/(2z)$.
They preserve all three eighth-power ratios by
the same explicit inverse-and-quotient root maps.
For $g_1$ leave $r$ unchanged.
For $g_2$ put $r'=rR/c$.
The polynomial identities of A.111 give
\[
 r'^4=-\frac{F(g_2c)}{F(g_2R)},\qquad
 f'=\frac{g_2T-g_2c}{g_2R-g_2c}
                  =f\frac RT,\qquad
 \Psi_{r'}-\Psi_r=\chi_{R/c},
        \qquad f=\frac{T-c}{R-c}.
\tag{A.116.5}
\]
The factor $R/T$ is an eighth power and its cyclic
algebra with an order-four character is zero.
The additional quadratic symbol in (A.116.5) is also
zero for an explicit reason.
For $D=c/R$ and $T/R=X^8$,
\[
 f=\frac{X^8-D}{1-D},\qquad
 X^8-D=N_{k(\sqrt D)/k}(X^4+\sqrt D),\qquad
 1-D=N_{k(\sqrt D)/k}(1+\sqrt D).
\tag{A.116.6}
\]
These identities include the split quadratic algebra.
Therefore $(f,\chi_{R/c})=0$, since $R/c$ and $D$
have the same squareclass.
Equations (A.116.5)--(A.116.6) prove equality of the
original and transformed Brauer classes.

There is one more exact operation that will be useful.
Choose the three actual eighth roots in (A.116.1), and
write $W$ for their product.
Then $F(T)/F(R)=W^8$.
Exchanging $R$ and $T$ sends
\[
             r\longmapsto r/W^2,\qquad
             f\longmapsto f^{-1},\qquad
             \Psi_{r/W^2}=\Psi_r.
\tag{A.116.6a}
\]
Consequently this exchange negates the displayed Brauer
class and preserves its vanishing.

\emph{The dyadic place.}
Now take $k=\Q_2$.
The two coordinate transformations place $e=v_2(R)$ in the range
$e\ge1$, by the same three valuation cases as A.111.

\emph{The exact valuation bounds for the eighth-power cover.}
Write $e_T=v_2(T)$.
The valuation congruences are now
\[
 v_2(T)\equiv e,\qquad v_2(T-1)\equiv0,\qquad
                  v_2(T-\tfrac12)\equiv-1\pmod8.
\tag{A.116.7}
\]
Negative $e_T$ contradict the last two congruences.
If $e_T=0$, both $e$ and $v_2(T-1)$ are positive
multiples of eight.
Then $(1-2T)/(1-2R)\equiv-1\pmod{32}$, whereas
an eighth-power dyadic unit is one modulo thirty-two.
Thus $e_T>0$.
The ratio $(1-T)/(1-R)$ also gives
$v_2(T-R)\ge5$.
The ratio $T/R$ gives $e_T\equiv e\pmod8$ and,
when $e_T=e$, the stronger bound $e+5$.
Consequently
\[
 \begin{array}{ll}
 1\le e\le4:&e_T=e,\quad v_2(T-R)\ge e+5,\\
 e\ge5:&e_T\ge5,\quad e_T\equiv e\pmod8.
 \end{array}
\tag{A.116.8}
\]
The unequal-depth cases follow by taking the
minimum of the two valuations, exactly as in A.111.

Write $F(z)=zG(z)/2$ with
$G(z)=(1-z)(1-2z)$.
If $k=v_2(c)\le0$, then
$v_2(f-1)=v_2(T-R)-k\ge5$, so $f$ is a fourth
power and (A.116.2) follows.
Assume henceforth that $k>0$.
All three $G$-values are units and (A.116.1) gives
\[
                         k-e=4v_2(r).
\tag{A.116.9}
\]

For $1\le e\le4$, if $k\ne e$ then
$v_2(R-c)\le e$ and (A.116.8) gives $f\in1+32\Z_2$.
If $k=e$, write $c=Rz$ with $z$ odd.
The possibility $v_2(z-1)\ge2$ is excluded:
\[
 G(Rz)-G(R)=R(z-1)(-3+2R(z+1))
\tag{A.116.10}
\]
is divisible by eight, and
$r^4=-zG(Rz)/G(R)\equiv-z\equiv3\pmod4$.
Thus $v_2(R-c)\le e+1$, giving $f\in1+16\Z_2$.
Every such unit is a fourth power, so its
cyclic algebra with $\Psi_r$ is zero.

For $e\ge5$ and $1\le k\le4$, (A.116.9) and
$e_T\equiv e\pmod8$ make
$\min(e,e_T)-k$ a positive multiple of four.
Hence again $f\in1+16\Z_2$, proving (A.116.2).

It remains that $e,k\ge5$.
Here $G(c)/G(R)\in1+32\Z_2$.
It has a fourth root $l\in1+8\Z_2$.
For clarity this follows from
\[
 (1+8w)^4=1+32(w+12w^2+64w^3+128w^4):
\tag{A.116.11}
\]
the polynomial in parentheses has derivative one
modulo two and hence takes every value of $\Z_2$.
The unit $l$ is itself a square.
Set $b=r/l$.
Then $b^4=-c/R$ and $\Psi_b=\Psi_r$.
With the actual eighth root $X$ of $T/R$,
\[
                         f=\frac{X^8+b^4}{1+b^4}.
\tag{A.116.12}
\]
Both values are the norms in (A.116.4).
Their ratio is a norm as well, which proves
(A.116.2) in this last case.

\emph{Every odd finite place.}
Let $k=\Q_p$ with $p$ odd.
The four punctures have distinct reductions in
$\mathbf P^1(\F_p)$.
Call a point good when its reduction avoids these
four reductions; otherwise specify the boundary
residue disc containing it.
The maps $g_1,g_2$ belong to $\operatorname{PGL}_2(\Z_p)$,
preserve good points, and permute the four boundary
discs transitively.
Every principal unit is an eighth power by Hensel's
lemma, since eight is a unit in $\Z_p$.

The character $\chi_{16}$ is unramified at $p$:
adjoining sixteenth roots of unity is etale over
$\Z_p$, as the derivative of $Z^{16}-1$ is a unit
at each root.
If $v_p(r)$ is even, the quadratic character
$\chi_r$ is unramified as well, by adjoining the
square root of its unit part.
For an unramified cyclic extension of degree
$m\mid4$, every unit and every element of valuation
divisible by eight is a norm.
Indeed the finite-field norm is onto, and the norm
on principal units is onto because every principal
unit already has an $m$th root in the base field;
the norm of the base uniformizer is its $m$th power.
These observations prove the exact criterion
\[
 v_p(r)\in2\Z,\qquad v_p(f)\in8\Z
           \quad\Longrightarrow\quad (f,\Psi_r)=0.
\tag{A.116.13}
\]

First suppose $R$ and $c$ lie in the same boundary
disc, and send that disc to zero.
Then $R,c\in p\Z_p$ and $G(c)/G(R)$ is a principal
unit.
Choose an actual eighth root $u$ of this ratio and
put $b=r/u^2$.
The source equation gives
\[
 b^4=-c/R,\qquad \Psi_b=\Psi_r,\qquad
 f=\frac{X^8+b^4}{1+b^4},\qquad X^8=T/R.
\tag{A.116.14}
\]
Thus the actual norm comparison (A.116.4) proves
vanishing, for every position of $T$.

If $R,T$ lie in the same boundary disc and $c$ lies
elsewhere, send the first disc to zero.
Now $v_p(R),v_p(T)>0$ and $v_p(c)\le0$, so
\[
                   f=\frac{1-T/c}{1-R/c}\in1+p\Z_p.
\tag{A.116.15}
\]
This is an eighth power, and again its symbol is zero.

If exactly one of $R,T$ is good, exchange them first
by (A.116.6a), so that $R$ is good.
Suppose henceforth that $R$ is good.
If $c$ is in a boundary disc, send that disc to zero.
If $T$ is also in that disc, exchanging $R,T$
reduces to (A.116.14).
Otherwise $v_p(c)>0$, $v_p(R)=0$, and $v_p(T)\le0$.
The exact identity
\[
                  \frac{f}{T/R}
                   =\frac{1-c/T}{1-c/R}\in1+p\Z_p
\tag{A.116.16}
\]
makes $f$ an eighth power.

It remains in this case that $R,c$ are both good.
Then $F(R),F(c)$ are units, so $r$ is a unit and
$\Psi_r$ is unramified.
If $R$ reduces to $c$, every
$D_i=(c-c_i)/(R-c_i)$ is a principal unit and hence
an eighth power.
If $T$ is good and reduces to $c$, every $D_i$ is
the product of the old eighth-power ratio
$(T-c_i)/(R-c_i)$ and a principal unit.
In either event the three-root fibre has a
$k$-rational root tuple.
Its actual product root and (A.116.3) imply
$K_r=k$, so the character itself is zero.
In all other events $v_p(f)=0$, except when $T$
lies in the disc of infinity.
There the old ratio $(T-1)/(R-1)$ gives
$v_p(T)\in8\Z$, and $v_p(f)=v_p(T)$.
Criterion (A.116.13) proves vanishing.

The only remaining configuration has $R,T$ in two
distinct boundary discs.
If $c$ shares either disc, (A.116.14), with an
exchange when needed, applies.
If $c$ lies in a third boundary disc, send it to
zero; both $R,T$ then have nonpositive valuation,
so (A.116.16) applies.
Finally suppose $c$ is good, and send the disc of
$R$ to zero.
Write $e=v_p(R)>0$.
If $T$ is in another finite boundary disc, it is a
unit, and $T/R\in k^{\times8}$ gives $e\in8\Z$.
If $T$ is in the disc of infinity, the ratio
$(T-1)/(R-1)$ first gives $v_p(T)\in8\Z$; the
ratio $T/R$ then again gives $e\in8\Z$.
Since $F(c)$ is a unit and $v_p(F(R))=e$, we obtain
\[
 v_p(r)=-e/4\in2\Z,\qquad
 v_p(f)=
 \begin{cases}
  0,&T\text{ in a finite boundary disc},\\
  v_p(T)\in8\Z,&T\text{ in the disc of infinity}.
 \end{cases}
\tag{A.116.17}
\]
Criterion (A.116.13) finishes the odd-place proof.
The cases just listed cover every reduction of
$R,T,c$.

\emph{The real place.}
Each of the three ratios in (A.116.1) is positive.
Consequently $R,T$ belong to the same connected
interval of $\R\setminus\{0,\tfrac12,1\}$.
The equation $-F(c)/F(R)=r^4>0$ implies that $c$
does not belong to that interval, since $F$ has
constant sign there.
Thus $(T-c)/(R-c)>0$.
A positive real number is a norm from $\C$;
every character of $G_{\R}$ has order at most two,
so the corresponding cyclic algebra splits.
This proves (A.116.2) also at infinity.
\end{proof}

The retained search
\path{code/harmonic_eighth_character_search.py}
computes the same characters on exact local eighth-power
points. Its output records 113953 trials and 13694
actual local points across twenty-four rational
negative-fourth-power sources, with no nonzero value.
The proof above, rather than this finite search,
establishes the all-place assertion for the stated family.
The separate exact certificate
\path{code/harmonic_eighth_dyadic_certificate.py},
with retained output, checks every finite cyclotomic
root-stabilizer case, the coordinate identities, the
octic norm resultant and the principal-unit fourth-root
polynomial used in this proof.

